\section{第三讲：轨道积分}\label{s:orb}
仿射 Springer 纤维在表示论中的意义，体现在它们与轨道积分的密切关系上。轨道积分是 $p$-进群调和分析中出现的某类积分。正如有限群的共轭类对理解其表示是基本的，轨道积分对理解 $p$-进群的表示也是基本的。在某些情形下，轨道积分可以解释为仿射 Springer 纤维上的点数。

\subsection{$p$-进群上的积分}\label{ss:int}

\subsubsection{设定} 设 $F$ 是局部非阿基米德域，即 $F$ 或是 $\QQ_{p}$ 的有限扩张，或是 $\FF_{p}\lr{t}$ 的有限扩张。则 $F$ 有离散赋值 $\val: F^{\times}\to \ZZ$，我们把它归一化为满射。设 $\calO_{F}$ 是 $F$ 的赋值环，$k$ 是剩余域。因此与前面几节不同，这里的 $k$ 是有限域。假定 $\chark$ 相对所讨论的群足够大。

\subsubsection{Haar 测度与积分} 设 $G$ 是 $F$ 上代数群。拓扑群 $G(F)$ 局部紧且完全不连通。它有右不变 Haar 测度 $\mu_{G}$，在相差一个标量的意义下唯一。对可测子集 $S\subset G(F)$，用 $\vol(S,\mu_{G})$ 记它在 $\mu_{G}$ 下的体积。固定紧开子群 $K_{0}\subset G(F)$，可以把 Haar 测度 $\mu_{G}$ 归一化，使 $K_{0}$ 的体积为 $1$。例如，若取 $G$ 在 $\calO_{F}$ 上的整模型 $\calG$，则可取 $K_{0}=\calG(\calO_{F})$。

借助 Haar 测度 $\mu_{G}$，可以对 $G(F)$ 上光滑（即局部常值）、紧支、复值的函数积分。把此函数空间记为 $\sS(G(F))$（其中 $\sS$ 取自 Schwarz）。对 $f\in\sS(G(F))$，积分
\begin{equation*}
\int_{G(F)}f\mu_{G}
\end{equation*}
可以如下计算。可以找到有限指数的子群 $K\subset K_{0}$，使 $f$ 右 $K$-不变，即对一切 $g\in G(F)$、$x\in K$ 有 $f(gx)=f(g)$（见习题 \ref{exIII:K inv}）。于是上述积分化为陪集 $G(F)/K$ 上的加权计数：
\begin{equation*}
\int_{G(F)}f\mu_{G}=\vol(K,\mu_{G})\sum_{[g]\in G(F)/K}f(g)=\frac{1}{[K_{0}:K]}\sum_{[g]\in G(F)/K}f(g).
\end{equation*}
容易验证：只要 $f$ 右 $K$-不变、且 $K$ 在 $K_{0}$ 中有限指数，上式右端就与 $K$ 的选取无关。

\subsubsection{变体} 设 $H\subset G$ 是定义在 $F$ 上的子群，并带有 Haar 测度 $\mu_{H}$。考虑函数 $f\in\sS(H(F)\bs G(F))$，即 $f$ 是 $G(F)$ 上左 $H(F)$-不变、局部常值、且支集模 $H(F)$ 紧的函数，则可以定义积分
\begin{equation}\label{int quot}
\int_{H(F)\bs G(F)}f\frac{\mu_{G}}{\mu_{H}}.
\end{equation}
此积分的算法如下。同样取有限指数的子群 $K\subset G(F)$，使 $f$ 右 $K$-不变。则积分 \eqref{int quot} 可以写成双陪集 $H(F)\bs G(F)/K$ 上的加权和：
\begin{equation*}
\int_{H(F)\bs G(F)}f\frac{\mu_{G}}{\mu_{H}}=\frac{1}{[K_{0}:K]}\sum_{[g]\in H(F)\bs G(F)/K}\frac{f(g)}{\vol(H(F)\cap gKg^{-1}, \mu_{H})}.
\end{equation*}

\subsection{轨道积分}

\subsubsection{轨道积分的定义} 沿用 \S\ref{ss:int} 的设定。把 $G$ 的 Lie 代数记为 $\frg(F)$，以强调它是 $F$ 上的向量空间。设 $\varphi\in\sS(\frg(F))$，$\gamma\in\frg(F)$。考虑映射 $G(F)\to \frg(F)$，$g\mapsto \Ad(g^{-1})\gamma$。则复合 $f: g\mapsto \varphi(\Ad(g^{-1})\gamma)$ 是 $G(F)$ 上的光滑函数。于是 $f$ 局部常值，在 $\gamma$ 于 $G(F)$ 中的中心化子 $G_{\gamma}(F)$ 下左不变，且支集模 $G_{\gamma}(F)$ 紧。

固定 $G(F)$ 上的 Haar 测度 $\mu_{G}$ 与 $G_{\gamma}(F)$ 上的 Haar 测度 $\mu_{G_{\gamma}}$。下面这个积分是 \eqref{int quot} 的特例（只是我们下面把积分变量 $g\in G_{\gamma}(F)\bs G(F)$ 明确写出，而 \eqref{int quot} 中没有）：
\begin{equation*}
O_{\gamma}(\varphi):=\int_{G_{\gamma}(F)\bs G(F)}\varphi(\Ad(g^{-1})\gamma)\frac{\mu_{G}}{\mu_{G_{\gamma}}}.
\end{equation*}
这类积分称为 Lie 代数 $\frg(F)$ 上的{\em 轨道积分}\zy{式中的 $\mu_{G}/\mu_{G_{\gamma}}$ 不是两个测度相除，而是 $G_{\gamma}(F)\bs G(F)$ 上由 $\mu_{G}$ 与 $\mu_{G_{\gamma}}$ 定出的商测度的记号，满足 $\int_{G}f\,\mu_{G}=\int_{G_{\gamma}\bs G}\big(\int_{G_{\gamma}}f(hg)\,\mu_{G_{\gamma}}\big)\,\mathrm{d}(\mu_{G}/\mu_{G_{\gamma}})$。来源：原文 \S3.1 的 \eqref{int quot} 与此处文脉。}。把 $\gamma$ 换成 $G(F)$ 中的元素、$\varphi$ 换成 $\sS(G(F))$ 中的元素，同样可以定义群 $G(F)$ 上的轨道积分。

\subsubsection{具体情形} 本节余下部分限制在如下情形。设 $G$ 是 $F$ 上分裂约化群。把相应的 Chevalley 群概形从 $\ZZ$ 基变换到 $\calO_{F}$，就固定了 $G$ 的一个整模型。以下把 $G$ 视为 $\calO_{F}$ 上的约化群概形。要求 $K_{0}=G(\calO_{F})$ 的体积为 $1$，以此归一化 $G(F)$ 上的 Haar 测度 $\mu_{G}$。

Lie 代数 $\frg(F)$ 含有来自 $\calO_{F}$ 上整模型的典范格 $\frg(\calO_{F})$。我们最感兴趣的是格 $\frg(\calO_{F})$ 的特征函数 $\varphi=\one_{\frg(\calO_{F})}$ 的轨道积分。

\subsubsection{$\gamma$ 的中心化子} 设 $\gamma$ 正则半单，于是其中心化子 $G_{\gamma}$ 是 $F$ 上环面。令 $F^{\ur}=F\wh{\otimes}_{k}\kbar$，它是剩余域代数闭的完备离散赋值域。仍用 $t$ 记 $F$ 的素元，它也是 $F^{\ur}$ 的素元。用 $t$，\S\ref{sss:L gamma} 中的构造给出嵌入 $\xcoch(G_{\gamma}\otimes_{F}F^{\ur})\incl  G_{\gamma}(F^{\ur})$，其像仍记为 $\L_{\gamma}$。此嵌入是 $\Gk$-等变的，故 $\L_{\gamma}$ 带有 $\Gk$ 的作用。把 $\L_{\gamma}$ 看成 $k$ 上的 étale 群概形，则记号 $\L_{\gamma}(k)$ 有意义；若把 $\L_{\gamma}$ 看成普通群，它就恰是 $\L_{\gamma}$ 中的 $\Gk$-不变量。于是 $\L_{\gamma}(k)\subset G_{\gamma}(F)$ 是离散余紧子群。

\subsubsection{$\GL_{n}$ 中的中心化子}\label{sss:GLn orb} 设 $G=\GL_{n}$，$\gamma$ 是 $\frg(F)$ 中的正则半单元，不要求它在 $F$ 上可对角化。假定 $\textup{char}(F)=0$ 或 $\chark>n$。如 \S\ref{sss:GLn centralizer}，$\gamma$ 的特征多项式 $P(x)=x^{n}+a_{1}x^{n-1}+\cdots+a_{n}\in F[x]$ 可分，故 $F$ 代数 $F[x]/(P(x))$ 同构于域之积 $F_{1}\times \cdots \times F_{m}$。有
\begin{equation*}
G_{\gamma}(F)\cong F^{\times}_{1}\times\cdots\times F^{\times}_{m}.
\end{equation*}
此时格 $\L_{\gamma}(k)\subset G_{\gamma}(F)=F^{\times}_{1}\times\cdots\times F^{\times}_{m}$ 由形如 $(t^{d_{1}},\cdots, t^{d_{m}})$（$(d_{1},\cdots, d_{m})\in\ZZ^{m}$）的元素组成。商 $\L_{\gamma}(k)\bs G_{\gamma}(F)$ 同构于 $\prod_{i=1}^{m}F^{\times}_{i}/t^{\ZZ}$。每个 $F^{\times}_{i}/t^{\ZZ}$ 落在正合列
\begin{equation}\label{Fi seq}
0\to \calO^{\times}_{F_{i}}\to F^{\times}_{i}/t^{\ZZ}\to \ZZ/e_{i}\ZZ\to 0
\end{equation}
中，其中 $e_{i}$ 是扩张 $F_{i}/F$ 的分歧指数，因此商 $\L_{\gamma}(k)\bs G_{\gamma}(F)$ 紧。

\subsubsection{用计数表示轨道积分}\label{sss:X gamma} 考虑 $G(F)/G(\calO_{F})$ 的如下子集
\begin{equation*}
X_{\gamma}:=\{[g]\in G(F)/G(\calO_{F})|\Ad(g^{-1})\gamma\in\frg(\calO_{F})\}.
\end{equation*}
这是仿射 Springer 纤维的集合论版本。

群 $G_{\gamma}(F)$ 按规则 $G_{\gamma}(F)\ni h:[g]\mapsto[hg]$ 作用在 $X_{\gamma}$ 上。对任意自由交换子群 $L\subset G_{\gamma}(F)$，它在 $G(F)/G(\calO_{F})$ 上的左平移作用自由（因为稳定子必有限），从而它在 $X_{\gamma}$ 上的作用也自由。

更一般地，对任意离散余紧子群 $L\subset G_{\gamma}(F)$，商群胚 $L\bs X_{\gamma}$ 是{\em 有限型}的，即有有限多个同构类，且每个对象的自同构群有限。对有限型群胚 $Y$，定义 $Y$ 的基数为
\begin{equation}\label{card groupoid}
\#Y:=\sum_{y\in \textup{Ob}(Y)/\cong}\frac{1}{\#\Aut(y)}.
\end{equation}

下一个引理直接来自定义，其证明留作习题 \ref{exIII:Ogamma}。

\begin{lemma}\label{l:Ogamma} 设 $\gamma$ 是 $\frg(F)$ 中的正则半单元，$L\subset G_{\gamma}(F)$ 是任意离散余紧子群。则
\begin{equation*}
O_{\gamma}(\one_{\frg(\calO_{F})})=\frac{1}{\vol(G_{\gamma}(F)/L,\mu_{G_{\gamma}})}\#\left(L\bs X_{\gamma}\right),
\end{equation*}
其中右端的基数按 \eqref{card groupoid} 解释。
\end{lemma}

\subsubsection{$G=\GL_{n}$ 的情形与分式理想}\label{sss:GLn orb formula} 沿用 \S\ref{sss:GLn orb} 的情形。在 $G(F)/G(\calO_{F})$ 与 $F^{n}$ 中 $\calO_{F}$-格的集合的等同（见 \eqref{Gr lattices}）下，有
\begin{equation*}
X_{\gamma}\cong\{\textup{格 }\L\subset F^{n}|\gamma \L\subset \L\}.
\end{equation*}
此双射把 $[g]\in X_{\gamma}$ 送到格 $\L=g\calO^{n}_{F}$。

再给出 $X_{\gamma}$ 的另一种解释。设 $P(x)=x^{n}+a_{1}x^{n-1}+\cdots+a_{n}$ 是 $\gamma$ 的特征多项式。令
\begin{equation*}
A=\calO_{F}[x]/(P(x))
\end{equation*}
为 $\frg(F)$ 中由 $\gamma$ 生成的交换 $\calO_{F}$-子代数。全分式环 $\Frac(A)$ 是次数为 $n$ 的有限 étale $F$-代数，而 $A$ 是其中的一个序。$A$ 在 $F^{n}$ 上的典范作用把 $F^{n}$ 实现为秩 $1$ 的自由 $\Frac(A)$-模。回忆{\em 分式 $A$-理想}指有限生成的 $A$-子模 $M\subset \Frac(A)$。若取元素 $e\in F^{n}$ 作为该 $\Frac(A)$-模结构的基，就得到双射
\begin{equation}\label{frac ideals}
\{\mbox{分式 $A$-理想}\}\bij X_{\gamma},
\end{equation}
它把 $M\subset \Frac(A)$ 送到 $\calO_{F}$-格 $M\cdot e\subset F^{n}$。

借助代数 $A$，有典范同构
\begin{equation*}
G_{\gamma}(F)\cong \Frac(A)^{\times}.
\end{equation*}
此同构使 $G_{\gamma}(F)$ 通过左平移在 $X_{\gamma}$ 上的作用，与 $\Frac(A)^{\times}$ 通过乘法在分式 $A$-理想集合上的作用相互缠绕。

当 $A$ 恰是 Dedekind 整环之积（即 $A$ 是 $\Frac(A)$ 中的极大序）时，所有分式 $A$-理想都是主理想；这等价于说 $G_{\gamma}(F)\cong \Frac(A)^{\times}$ 在 $X_{\gamma}$ 上的作用可迁。一般情形下，主分式理想构成 $G_{\gamma}(F)$ 下的齐性空间 $\Frac(A)^{\times}/A^{\times}$；一般情形下计数 $X_{\gamma}$ 的困难来自环 $A$ 的奇异性。

把 $G_{\gamma}(F)\cong\prod_{i}F^{\times}_{i}$ 上的 Haar 测度 $\mu_{G_{\gamma}}$ 归一化，使 $\prod_{i}\calO^{\times}_{F_{i}}$ 的体积为 $1$。取素元 $t_{i}\in F_{i}$，可作出格 $L_{0}=t^{\ZZ}_{1}\times \cdots\times t^{\ZZ}_{m}\subset G_{\gamma}(F)$。此时 $G_{\gamma}(F)/L_{0}\cong \prod_{i}\calO_{F_{i}}^{\times}$ 的体积为 $1$。因此引理 \ref{l:Ogamma} 给出
\begin{equation}\label{GLn L0}
O_{\gamma}(\one_{\frg(\calO_{F})})=\#(L_{0}\bs X_{\gamma}).
\end{equation}
借助 \eqref{frac ideals}，可以把 \eqref{GLn L0} 解释为：$O_{\gamma}(\one_{\frg(\calO_{F})})$ 等于分式 $A$-理想在乘以 $t_{i}$ 的各次幂之下的等价类的个数。
